\section{The 235-state fold: two meanings recovered from a preceding bit}
\label{sec:fold235}

Let $M$ denote the 236-state basis after the seven independently certified
source reassociations, and let $N$ denote the resulting 235-state machine.
All state labels in this section are inherited labels.  Both machines have
initial state $95$ and a separate halting outcome $254$.  A configuration is
$(q,h,T)$, where $h\in\mathbb Z$ and $T:\mathbb Z\to\{0,1\}$.
The source reassociations are a separate part of the construction: this
section proves an unconditional local equivalence between their exact
236-state output and the 235-state table.

\subsection{Exact edit and incoming invariant}

The four special old transition rows are
\[
\begin{array}{c|cc}
\text{Old state}&\text{Read }0&\text{Read }1\\\hline
74&0R114&1R233\\
193&0R190&1R213.
\end{array}
\]
Delete state $193$, redirect its incoming rows $163/1$ and $229/1$ to $74$,
and replace the two rows of state $74$ by
\[
74/0=0L138,\qquad 74/1=1L116.
\]
Every other retained row is unchanged.  Define
\[
\phi(q)=\begin{cases}74,&q=193,\\q,&q\ne193,\end{cases}
\qquad
F(q,h,T)=(\phi(q),h,T).
\]
The map fixes both the initial state and the halt outcome.  There are exactly
$236$ old working states with $472$ rows and $235$ new working states with
$470$ rows.  No extra initial tape word, tape symbol, or intermediate control
state is supplied.

The old machine satisfies the invariant
\[
I(q,h,T):\quad
(q=74\Longrightarrow T(h-1)=0)\ \land\
(q=193\Longrightarrow T(h-1)=1).
\]
Indeed, the complete old incoming-edge inventory is
\[
\begin{array}{c|c|c}
\text{Target}&\text{Incoming rows}&\text{Written bit and move}\\\hline
74&140/0&0,R\\
193&163/1,\ 229/1&1,R.
\end{array}
\]
Each incoming transition therefore puts the required bit immediately left
of the new head.  The initial state is neither special state, so $I$ holds
initially for every tape and every head position.  Induction on the old
step count establishes $I$ for every reachable old configuration.  The
invariant is asserted at old configurations and their simulation
checkpoints; the reused helper states may have additional meanings between
those checkpoints.

\subsection{Four three-step identities on open tapes}

Put the initial head at coordinate zero for notation, write $b=T(0)$, and
write $\ell=T(-1)$.  For the four old cases the complete new paths are
\[
\begin{array}{c|c|c|c}
\text{Old state}&b&\ell&\text{New control path}\\\hline
74&0&0&74\to138\to210\to114\\
74&1&0&74\to116\to246\to233\\
193&0&1&74\to138\to211\to190\\
193&1&1&74\to116\to134\to213.
\end{array}
\]
Each displayed path has exactly three transitions.  Their head positions
are $0,-1,0,1$.  Their read bits are $b,\ell,b$, respectively, and each
transition writes back the bit it reads.  Thus every cell is preserved,
including all unspecified exterior cells, and the final head displacement
is $+1$.  The endpoint states $114,233,190,213$ are precisely the endpoints
of the corresponding single old transitions.  None is an interior halt.

Equivalently, the two new rows first choose the helper according to the
current bit.  Helper $138$ or $116$ then reads the preceding bit, recovering
which old control state was intended.  One final existing decoder row
completes that old transition.  The source reassociations supply the
necessary four continuations without introducing a helper state.

These identities quantify over every surrounding binary tape.  The finite
footprint $\{-1,0\}$ is a consequence of the displayed paths, not a bound on
possible tape contents or future computation.  It suffices to check the
four prescribed assignments to this footprint because the invariant and
the current read bit fix both accessed cells.  The remaining $468$ old rows
commute with $F$ in exactly one new step, including the two incoming
redirections implementing $\phi$.

\subsection{Halting equivalence and exact clock identity}

\paragraph{Theorem.}
For every binary tape $T$ and every head position $h$,
\[
\operatorname{Halts}(M,95,h,T)
\quad\Longleftrightarrow\quad
\operatorname{Halts}(N,95,h,T).
\]

\paragraph{Proof.}
Let $C_t$ be the old run while it is defined.  Set $\tau(0)=0$ and
\[
\tau(t+1)-\tau(t)=
\begin{cases}
3,&\text{the control state of }C_t\text{ is }74\text{ or }193,\\
1,&\text{otherwise}.
\end{cases}
\]
The inductive incoming invariant and the four block identities prove
$N^{\tau(t)}(F(C_0))=F(C_t)$ by induction on $t$.  An old halt is therefore
a new halt.  Conversely, an infinite old run supplies infinitely many
positive new blocks; their synchronization times satisfy $\tau(t)\ge t$,
and no block contains an interior halt.  Their union covers the entire new
run, which consequently cannot halt.  This establishes both directions.
\hfill$\square$

If $V_q(t)$ counts old visits to state $q$ among the first $t$ transitions,
the proof gives the exact identity
\[
\tau(t)=t+2V_{74}(t)+2V_{193}(t),\qquad t\le\tau(t)\le3t.
\]
The theorem uses only finite table calculations and induction and is
formalizable in ZFC.  Composing it with the separately verified source
reassociation, initialization, runtime simulation, and inherited logical
calibration transfers the desired equivalence with
$\operatorname{Con}(\mathrm{ZFC})$.  The local theorem itself requires no
consistency hypothesis and does not assert an unconditional nonhalting
result.

\subsection{Independent executable verification}

The delivered \texttt{check\_fold235.py} has no imports from either discovery
scanner or the generic symbolic checker.  It checks the exact input-byte
hashes, complete state and row inventories, all incoming edges, all $468$
ordinary row identities, and all four three-step blocks.  Its tape map
contains only the two fixed cells; an attempted exterior read is rejected.
Every write is checked against the current value, and every complete
state/head/tape endpoint is compared with the old transition.  Its receipt
records the full traces.

A second route, the inherited generic open-tape verifier, independently
rederives the bit invariant and proves all $472$ old rows, producing $468$
one-step blocks and four three-step blocks.  The two routes support the
ordinary mathematical proof above.  They do not constitute a newly
compiled proof-assistant theorem for the complete 235-state machine.
